\section{Discussion}
\label{sec:discussion}

\paragraph{Potential applications.}
Our findings suggest a way to evaluate the reuse of contextual skills.
Consider a Web agent that follows the same instructions to extract
specified information from different pages. The instructions stay the
same, while the page content and the user's request change. Reusing the
states formed from those instructions would need to preserve the agent's
ability to answer the new request. The separation observed here motivates
testing that ability directly, alongside the benefit of providing the
original instructions for each request.

An evaluation can first compare a receiving execution that uses the stored
states with an execution that processes the correct skill for that request.
This measures how much task performance survives reuse. It can then hold
the copied skill information fixed and compare retaining the receiving
execution's existing caches with recomputing the later prompt. That
comparison checks whether the copied information works with the history
already stored by the receiving execution. The complete-source controls
motivate this second check: matching all skill-position KV still leaves
performance sensitive to the state of other positions. When task inputs
change, the amount of prompt processing that can be reused therefore
needs to be assessed together with the skill representation.

Continued access to the skill is a further part of the evaluation.
A method that preserves accuracy after copying states may still require
attention to the skill positions. Comparing runs with and without this
access distinguishes reuse that depends on continued skill reading from
reuse that does not. This distinction matters when the intended
benefit is to reduce repeated instruction processing. Studies in
interactive Web settings can evaluate task performance, latency and
memory use together, counting both the copied cache and any later prompt
computation required to make use of it.

\paragraph{Future work.}
A next step is to investigate which questions benefit from cache
replacement. The mathematical analysis connects remaining cache
differences to output stability under explicit assumptions. Measuring
these differences, their influence on token predictions and the donor's
prediction margins could test whether this framework helps distinguish
successful from unsuccessful repairs. The complete-source controls further motivate testing which parts of
the receiving history must be recomputed, and whether the required budget
can be predicted on held-out skills. The heterogeneous depth effects suggest that a recomputation rule should
be tested across model/task cells rather than inferred from one depth curve. Separating attention cuts during prompt
processing from those during answer generation would clarify when continued
skill access matters. Finally,
multi-step Web workflows would assess how far the observed separation
extends beyond the question-answering settings studied here.
